\documentclass[11pt]{article}
\usepackage[T1]{fontenc}
\usepackage{lmodern}
\usepackage[a4paper,margin=28mm]{geometry}
\usepackage{amsmath,amssymb,amsthm,booktabs,microtype}
\usepackage{url}
\usepackage[colorlinks=true,linkcolor=blue,citecolor=blue,urlcolor=blue]{hyperref}
\usepackage{listings}
\hypersetup{pdftitle={Sharp local-dimension jumps and degree-jump growth in the partition graph},
  pdfauthor={Xiaojian Zeng},pdfsubject={Quantitative edgewise bounds for the integer partition graph}}



\newtheorem{theorem}{Theorem}
\newtheorem{lemma}[theorem]{Lemma}
\newtheorem{proposition}[theorem]{Proposition}
\newcommand{\supp}{\operatorname{supp}}
\newcommand{\one}{\mathbf{1}}
\newcommand{\HH}{\mathsf{H}}
\title{Sharp local-dimension jumps and degree-jump growth\\
in the partition graph}
\author{Xiaojian Zeng\\[3pt]
  \small PBC School of Finance, Tsinghua University, 100190\\
  \small \href{mailto:xiaojianzeng@foxmail.com}{xiaojianzeng@foxmail.com}}
\date{}

\begin{document}
\maketitle
\begin{abstract}
For the graph of integer partitions joined by a transfer of one unit between
rows, we determine the maximum change in local clique dimension along an
edge. It is zero for $n=2,3$, one for $4\leq n\leq10$ and $n=12,13$,
and two for $n=11$ and every $n\geq14$. We also characterize every edge
with dimension jump two by its support and multiplicities. If
$T_s=s(s+1)/2\leq n<T_{s+1}$ and $n\geq10$, the maximum degree jump
satisfies $4s-8\leq D_d(n)\leq4s+2$, so
$D_d(n)=4\sqrt{2n}+O(1)$. The arguments combine established vertex
formulas with restrictions on adjacent support patterns and explicit
witnesses for every required $n$. Only the finitely many small dimension
exceptions require exhaustive computation; a complete verifier is included.
These results answer the quantitative growth questions, while the original
problem's additional question about localization remains unresolved.
\end{abstract}

\section{The graph and the quantitative results}

Let $G_n$, for $n\geq2$, be the simple undirected graph whose vertices are
the integer partitions of $n$. A move selects a donor row of size $p\geq1$
and a different recipient row of size $q\geq1$, or a new row represented
by $q=0$. Replace these sizes by $p-1$ and $q+1$, discard zero parts,
and sort. The resulting partition must differ from the starting partition.
In particular, $p=q>0$ requires two distinct rows; repeated descriptions of
the same move create only one edge. This is the row-transfer convention of
Lyudogovskiy~\cite[Section~2]{LyudogovskiyMorphology}.

For a vertex $\lambda$, let $d(\lambda)$ be its degree and let
\[
 \delta(\lambda)=\max\{|C|-1:C\text{ is a clique of }G_n,
                                      \ \lambda\in C\}.
\]
Equivalently, $\delta(\lambda)$ is the clique number of the graph induced
by the neighbors of $\lambda$. Set
\[
 D_d(n)=\max_{\lambda\sim\mu}|d(\mu)-d(\lambda)|,
 \qquad
 D_\delta(n)=\max_{\lambda\sim\mu}|\delta(\mu)-\delta(\lambda)|.
\]
Lyudogovskiy~\cite[Problems~8.1--8.2]{LyudogovskiyJumps} asks for the
growth of these quantities and, separately, for the location of large
dimension jumps. Uniform boundedness of dimension jumps was also asked
in~\cite[Problem~6.4(iii)]{LyudogovskiyStratification}.

Write $r(\lambda)=|\supp(\lambda)|$ for the number of distinct positive
part sizes. Denote the staircase partition $(r,r-1,\ldots,1)$ by $S_r$.
We say that $\HH(\lambda)$ holds if a size $x\geq2$ occurs at least twice
in $\lambda$ and $x-1$ is absent.

\begin{theorem}\label{thm:dimension}
For every integer $n\geq2$,
\begin{equation}\label{eq:exact-dimension}
 D_\delta(n)=
 \begin{cases}
 0,& n=2,3,\\
 1,& 4\leq n\leq10\text{ or }n=12,13,\\
 2,& n=11\text{ or }n\geq14.
 \end{cases}
\end{equation}
Orient an edge $\lambda\to\mu$ so that $r(\mu)\geq r(\lambda)$ and put
$k=r(\mu)-r(\lambda)$. Its absolute dimension jump is two if and only if
one of the following holds:
\begin{enumerate}
\item $k=1$, $\lambda$ is neither a staircase nor an $\HH$ partition,
      and $\HH(\mu)$ holds;
\item $k=2$ and both $\HH(\lambda)$ and $\HH(\mu)$ hold.
\end{enumerate}
In either case $\delta(\mu)-\delta(\lambda)=2$.
\end{theorem}

\begin{theorem}\label{thm:degree}
Let $n\geq10$, $T_s=s(s+1)/2$, and choose the unique integer $s$ with
$T_s\leq n<T_{s+1}$. Then
\begin{equation}\label{eq:degree-bounds}
 4s-8\leq D_d(n)\leq4s+2.
\end{equation}
Consequently $D_d(n)=4\sqrt{2n}+O(1)$, with an absolute implied constant.
The lower bound has an explicit witness for every $n\geq10$.
\end{theorem}

The prior vertex formulas give degree in terms of support, gaps and
multiplicities~\cite[Theorem~2.2]{LyudogovskiyDegree}, and local dimension
in terms of corner capacities~\cite[Theorem~7.1]{LyudogovskiyMorphology}.
Support birth and death accounting, including the bound
$|r(\mu)-r(\lambda)|\leq2$, is established
in~\cite[Theorem~3.4 and Corollary~3.5]{LyudogovskiySupport}.
These inputs alone yield the coarse dimension-jump bound four.
The additional step here is to restrict which endpoint types can be
adjacent: this gives the sharp bound two and the equality criterion.
For degree, controlling the change in the correction terms and constructing
edges of almost maximal support difference yields the leading coefficient.
The vertex maxima studied in~\cite{LyudogovskiyDegree} are distinct from
the adjacent-pair maximum considered here.

Our literature comparison is bounded: no equivalent sharp edgewise
statement was located in the cited partition-graph preprints inspected on
14 September 2026. This is not an exhaustive priority assessment. The
conclusions below concern quantitative extrema and a local equality
criterion; the original localization question remains unresolved.

\section{Vertex formulas and restrictions on support}

Write the distinct sizes as $s_1>\cdots>s_r>0$, with multiplicities $m_i$,
and put $s_{r+1}=0$. Define
\[
 A=\#\{i:s_i-s_{i+1}>1\},\qquad B=\#\{i:m_i>1\},
\]
and set $a_i=\one_{\{m_i=1\}}$, $b_i=\one_{\{s_i-s_{i+1}=1\}}$,
with $a_{r+1}=b_0=0$. The established formulas are
\begin{align}
 d(\lambda)&=r(r-1)+A+B,\label{eq:vertex-degree}\\
 \delta(\lambda)&=\max\left\{
  \max_{1\leq i\leq r}(r+1-a_i-b_i),
  \max_{1\leq j\leq r+1}(r-a_j-b_{j-1})\right\}.
  \label{eq:vertex-dimension}
\end{align}
Equation~\eqref{eq:vertex-degree} is
\cite[Theorem~2.2]{LyudogovskiyDegree};
equation~\eqref{eq:vertex-dimension} is
\cite[Theorem~7.1]{LyudogovskiyMorphology}, also reproduced with these
boundary conventions in~\cite[Background Proposition~3.1]{LyudogovskiySynthesis}.
The two families in~\eqref{eq:vertex-dimension} are the row and column
capacities of the local admissibility graph.

Put $\varepsilon(\lambda)=\delta(\lambda)-r(\lambda)$. An immediate
reformulation of the attributed dimension formula is
\begin{equation}\label{eq:three-types}
 \varepsilon(\lambda)=
 \begin{cases}
 -1,&\lambda=S_r,\\
 1,&\HH(\lambda),\\
 0,&\text{otherwise}.
 \end{cases}
\end{equation}
Indeed, every column capacity is at most $r$. A row capacity equals $r+1$
exactly when $a_i=b_i=0$, which is precisely $\HH$. If $\HH$ fails,
all rows have capacity at most $r$; every row has capacity $r-1$ exactly
when all $a_i=b_i=1$, giving $S_r$. For $S_r$, the first and last column
capacities are $r-1$, and any other column has capacity $r-2$. Otherwise
some row has capacity $r$. This reasoning includes $r=1$.

Only donor and recipient sizes can disappear in a move, and only their
output sizes can appear. Omitting zero therefore gives
\begin{equation}\label{eq:support-bound}
 |r(\mu)-r(\lambda)|\leq2.
\end{equation}
This also gives a direct proof of the known support bound. We need two
further restrictions.

\begin{lemma}\label{lem:staircase}
Every neighbor of $S_r$ has strictly smaller support.
\end{lemma}
\begin{proof}
The partition $S_1$ has no neighbor. At $S_r$ all sizes occur once, so
equal positive donor and recipient sizes are impossible. A move with
$p=q+1$ is an identity and is excluded. If $q=0$, then $p\geq2$;
size $p$ disappears and both outputs $p-1$ and $1$ were already present.
If $q>0$, both distinct sizes $p,q$ disappear: an output could restore
one of them only in the excluded case $p=q+1$. The only possible size
outside $\{1,\ldots,r\}$ is $r+1$. Thus at least two sizes are lost
and at most one is gained. Coincident output sizes cannot increase the
number of new sizes.
\end{proof}

\begin{lemma}\label{lem:two-births}
If a move increases support by two, its starting partition satisfies $\HH$.
\end{lemma}
\begin{proof}
There are at most two new sizes, so an increase of two requires two
distinct positive new sizes and no disappearance. In particular $p>1$,
$p-1$ was absent, and size $p$ survives. Since the move removes a row of
size $p$ and neither output was previously present, another row of size
$p$ must have existed. Thus $p$ was repeated and its predecessor was
absent. If $p=q$, survival requires at least three copies, which gives the
same conclusion. A unit donor or coincident outputs cannot produce two
distinct positive new sizes.
\end{proof}

\section{Dimension jumps}

\begin{proof}[Proof of the universal bound and equality criterion in
Theorem~\ref{thm:dimension}]
Orient the edge toward nondecreasing support, and write
$k=r(\mu)-r(\lambda)$, $e=\varepsilon(\lambda)$ and
$f=\varepsilon(\mu)$. Then
\begin{equation}\label{eq:signed-jump}
 \delta(\mu)-\delta(\lambda)=k+f-e.
\end{equation}
By~\eqref{eq:support-bound}, $k\in\{0,1,2\}$.
If $k=0$, Lemma~\ref{lem:staircase} excludes staircases at both endpoints,
so $e,f\in\{0,1\}$ and the difference lies in $[-1,1]$.
If $k=1$, the same lemma excludes a staircase at the starting endpoint,
so $e\in\{0,1\}$ and $f\in\{-1,0,1\}$; the difference lies in $[-1,2]$.
If $k=2$, Lemma~\ref{lem:two-births} gives $e=1$, and the difference
is $1+f\in[0,2]$.
Hence its absolute value is at most two. Equality is possible exactly
when $k=1,e=0,f=1$, or $k=2,e=f=1$. Substitution also proves sufficiency,
and~\eqref{eq:three-types} identifies these conditions with the stated
endpoint criterion. Ties in the support orientation are harmless.
\end{proof}

\begin{proof}[Proof of the attainment set in Theorem~\ref{thm:dimension}]
At $n=11$ and $n=14$, respectively, use
\[
 (4,4,3)\sim(4,4,2,1),\qquad
 (5,5,4)\sim(5,5,3,1).
\]
Both edges split one row into a predecessor row and a new unit row.
Their endpoint dimensions are $2$ and $4$ by~\eqref{eq:three-types}.
For every $n\geq15$, set $h=n-11\geq4$ and sort the multisets in
\begin{equation}\label{eq:dimension-family}
 (h,4,4,3)\sim(h,4,4,2,1).
\end{equation}
If $h=4$, the supports have sizes two and three; the starting vertex
is neither a staircase nor an $\HH$ partition, while the final vertex
satisfies $\HH$ at size four. Their dimensions are $2$ and $4$.
For $h\geq5$, the supports have sizes three and four. The only repeated
size at the starting vertex is four, whose predecessor three is present;
at the final vertex that predecessor is absent. Their dimensions are
$3$ and $5$. Thus every asserted value two is attained.

For $n\geq4$, the edge $(n)\sim(n-1,1)$ has dimensions $1$ and $2$.
At $n=2$ the graph is a single edge, and at $n=3$ it is a three-vertex
path; all local dimensions are one. The remaining upper bounds are the
finite calculations in Table~\ref{tab:small}. A complete integer-only
verifier appears in Appendix~\ref{app:verifier}. Together with the
universal bound and~\eqref{eq:dimension-family}, these calculations prove
\eqref{eq:exact-dimension}.
\end{proof}

\begin{table}[htbp]
\centering
\caption{Exhaustive graph and clique calculations for the finite exceptions.
Every vertex and every edge is included.}\label{tab:small}
\begin{tabular}{rrrr}
\toprule
$n$ & $|V(G_n)|$ & $|E(G_n)|$ & $D_\delta(n)$\\
\midrule
2&2&1&0\\
3&3&2&0\\
4&5&5&1\\
5&7&9&1\\
6&11&17&1\\
7&15&28&1\\
8&22&47&1\\
9&30&73&1\\
10&42&114&1\\
11&56&170&2\\
12&77&253&1\\
13&101&365&1\\
\bottomrule
\end{tabular}
\end{table}

\section{Degree jumps}

\begin{lemma}\label{lem:correction}
For every edge $\lambda\sim\mu$,
\begin{equation}\label{eq:correction}
 \left|d(\mu)-d(\lambda)
 -\bigl(r(\mu)(r(\mu)-1)-r(\lambda)(r(\lambda)-1)\bigr)\right|\leq8.
\end{equation}
\end{lemma}
\begin{proof}
View $A$ as a function of the support, including the gap to zero.
Inserting a size $x$ between consecutive augmented sizes $a>b\geq0$
replaces $\one_{\{a-b>1\}}$ by
\[
 \one_{\{a-x>1\}}+\one_{\{x-b>1\}}.
\]
Since both smaller gaps are positive integers, the change has absolute
value at most one. Insertion above the maximum adds either zero or one;
the same holds when inserting into an empty support. Deletion reverses
an insertion. A move changes membership at at most four positive sizes,
so $|A(\mu)-A(\lambda)|\leq4$.

Likewise $B=\sum_{x\geq1}\one_{\{m_x\geq2\}}$ changes at at most four
sizes, and each indicator changes by at most one, even when donor and
recipient sizes coincide. Thus $|B(\mu)-B(\lambda)|\leq4$.
Equation~\eqref{eq:vertex-degree} now gives~\eqref{eq:correction}.
\end{proof}

\begin{proof}[Proof of Theorem~\ref{thm:degree}]
The distinct positive sizes of a partition with support $r$ sum to at
least $T_r$, so $r\leq s$. Because $s\geq4$ and the endpoint supports
differ by at most two, the largest difference in the quadratic terms is
\[
 s(s-1)-(s-2)(s-3)=4s-6.
\]
Lemma~\ref{lem:correction} proves $D_d(n)\leq4s+2$.

For the lower bound, write $n=T_s+q$ with $0\leq q\leq s$, and let
\begin{align*}
 \mu&=(s+q,s-1,\ldots,1),\\
 \lambda&=(s+q,s-1,\ldots,4,3,3).
\end{align*}
The range $s-1,\ldots,4$ in the second line is empty when $s=4$.
Both partitions sum to $n$, and transferring the unit row of $\mu$ into
its row of size two gives $\lambda$. At $\mu$,
\[
 (r,A,B)=(s,\one_{\{q>0\}},0),
\]
whereas at $\lambda$,
\[
 (r,A,B)=(s-2,1+\one_{\{q>0\}},1).
\]
The extra gap is the bottom gap to zero; the sole repeated size is three.
Consequently
\[
 d(\mu)-d(\lambda)=s(s-1)-(s-2)(s-3)-2=4s-8.
\]
This works for every $q$ in its range, proving~\eqref{eq:degree-bounds}.
Finally, the triangular inequalities imply
$s<\sqrt{2n}<s+3/2$. Thus replacing $s$ by $\sqrt{2n}$ changes the
bounds only by an absolute constant, giving the asserted asymptotic.
\end{proof}

\section{Computation, scope and assistance}

The finite part of the proof examines all 371 vertices and 1,084 edges
for $2\leq n\leq13$. The verifier in the appendix computes actual
neighborhood clique numbers without either vertex formula. The first clique
found supplies a lower certificate, and exhaustive rejection of all larger
subsets supplies the matching upper certificate. There is no symmetry
quotient, floating-point comparison or unproved search cutoff.

For a methodological comparison, Zeng, Liu, Ratnavelu and
Ong~\cite[Section~4, Lemma~3 and Proposition~1]{ZengFiniteCertificates}
use integer compatibility predicates and two exact clique implementations
for a finite library of normalized ternary codes. Here the finite graph is
small enough for literal subset enumeration. A separately written checker
also reconstructed the partition graphs by coefficient-wise partition
generation and literal row transfers, using a different clique search;
its checks through $n=24$ agree with the stated formulas and bounds.
These additional checks are finite corroboration. The infinite conclusions
rest on the support arguments and witness families above.

The equality criterion describes individual maximal dimension transitions
from endpoint multiplicities. It does not determine their spatial or
statistical concentration relative to simplex layers or other regions.
Accordingly, the original localization question remains unresolved.
We do not determine $D_d(n)$ exactly for all $n$, or claim sharpness of
the additive constants in~\eqref{eq:degree-bounds}.

\paragraph{Computational and AI assistance.}
OpenAI language-model assistance was used in developing the arguments,
drafting verification programs, comparing sources and preparing the text.
A separate model context inspected the mathematical arguments and ran a
separately written exact verifier. These checks are neither external peer
review nor proof-assistant certification. The computations use Python
integers and finite sets; exhaustive computation is used in the proof only
for the explicitly stated finite exceptions.

\appendix
\section{Complete verifier for the finite exceptions}\label{app:verifier}

Save the following as \texttt{verify\_small.py} and run
\texttt{python3 verify\_small.py}. Only the Python standard library is
needed. The program prints the four columns of Table~\ref{tab:small}.

The partition recursion chooses each next part no larger than its
predecessor, so every nonincreasing positive partition occurs exactly once.
The neighbor routine tries every donor row and every different recipient
row, including a new zero row, and deduplicates the results. Therefore it
constructs precisely $G_n$. The clique routine tests neighborhood subsets
in decreasing cardinality and returns at the first clique. Every larger
subset has been rejected by a missing edge, so the returned size is the
actual $\delta(\lambda)$. Taking the largest edge difference then gives
the displayed maximum. The assertions check closure and symmetry and
compare the computed rows with the stated table; they do not determine
any adjacency or clique value.

\begin{lstlisting}[language=Python,basicstyle=\ttfamily\footnotesize,columns=fullflexible,keepspaces=true,showstringspaces=false,breaklines=true,emptylines=1,frame=single,framesep=5pt]
"""Reproduce the finite dimension-jump table; Python standard library only."""
from itertools import combinations


def partitions(total, cap):
    if total == 0:
        yield ()
    else:
        for a in range(min(total, cap), 0, -1):
            for tail in partitions(total - a, a):
                yield (a,) + tail


def neighbors(v):
    out = set()
    for i in range(len(v)):
        for j in range(len(v) + 1):
            if i == j:
                continue
            w = list(v) + [0]
            w[i] -= 1
            w[j] += 1
            w = tuple(sorted((x for x in w if x), reverse=True))
            if w != v:
                out.add(w)
    return out


def clique_number(vertices, adj):
    vertices = sorted(vertices)
    for size in range(len(vertices), -1, -1):
        for c in combinations(vertices, size):
            if all(w in adj[v] for v, w in combinations(c, 2)):
                return size
    raise AssertionError("The empty set is a clique")


expected = [(2, 1, 0), (3, 2, 0), (5, 5, 1), (7, 9, 1),
            (11, 17, 1), (15, 28, 1), (22, 47, 1), (30, 73, 1),
            (42, 114, 1), (56, 170, 2), (77, 253, 1), (101, 365, 1)]
print("n vertices edges D_delta")
for n in range(2, 14):
    vs = list(partitions(n, n))
    assert len(vs) == len(set(vs))
    adj = {v: neighbors(v) for v in vs}
    assert all(adj[v] <= set(vs) for v in vs)
    assert all(v in adj[w] for v in vs for w in adj[v])
    dims = {v: clique_number(adj[v], adj) for v in vs}
    edges = [(v, w) for v in vs for w in adj[v] if v < w]
    jump = max(abs(dims[v] - dims[w]) for v, w in edges)
    row = (len(vs), len(edges), jump)
    assert row == expected[n - 2], (n, row)
    print(n, *row)
\end{lstlisting}

\bibliographystyle{plainurl}
\bibliography{references}
\end{document}
